\documentclass[11pt,a4paper]{article}

\usepackage[utf8]{inputenc}
\usepackage[T1]{fontenc}
\usepackage{lmodern}
\usepackage[margin=2.2cm]{geometry}
\usepackage{amsmath,amssymb}
\usepackage{booktabs}
\usepackage{listings}
\usepackage{xcolor}
\usepackage{hyperref}

\definecolor{codegray}{rgb}{0.45,0.45,0.45}
\definecolor{backcolour}{rgb}{0.96,0.96,0.96}

\lstdefinestyle{codeblock}{
    backgroundcolor=\color{backcolour},
    basicstyle=\ttfamily\small,
    breaklines=true,
    numbers=none,
    frame=single,
    rulecolor=\color{black!15}
}

\lstdefinestyle{terminaloutput}{
    backgroundcolor=\color{backcolour},
    basicstyle=\ttfamily\footnotesize,
    breaklines=true,
    numbers=none,
    frame=single,
    rulecolor=\color{black!15}
}

\lstset{style=codeblock}

\title{\textbf{Assignment 3: Adversarial Search}}
\author{Artificial Intelligence (AI 2) --- OsloMet}
\date{October 2026}

\begin{document}

\maketitle

\section{The Halving Game}

The Halving game begins with $N = 5$. Players alternate between decrementing ($N - 1$) and halving ($\lfloor N/2 \rfloor$). The player whose turn starts with $N = 0$ wins.

\subsection{Output}
\begin{lstlisting}[style=terminaloutput]
The number is 5 and it is P1's turn
P1's action: --
The number is 4 and it is P2's turn
P2's action: --
The number is 3 and it is P1's turn
P1's action: /2
The number is 1 and it is P2's turn
P2's action: --
The number is 0 and P1 won
\end{lstlisting}

\subsection{Discussion}
$N = 5$ is a forced win for Player 1. P1 moves to $N = 4$, leaving P2 in a losing position (both $4 - 1 = 3$ and $4 // 2 = 2$ lead to positions P1 can win from). P2 decrements to 3, P1 halves to 1, and P2 is forced to decrement to 0. At $N = 0$ it is P1's turn, so P1 wins.

\section{The Bucket Game}

Player 1 (MAX) chooses bucket $A = \{-50, 50\}$, $B = \{3, 1\}$, or $C = \{-5, 15\}$. Player 2 (MIN) then chooses one of the values inside that bucket.

\subsection{Output}
\begin{lstlisting}[style=terminaloutput]
The state is (0, ['A', 'B', 'C']) and it is P1's turn
P1's action: B
The state is (1, [3, 1]) and it is P2's turn
P2's action: 1
The state is (0, [1]) and P1's utility is 1
\end{lstlisting}

\subsection{Discussion}
Minimax evaluates the worst-case payoff of each bucket:
\begin{itemize}
    \item Bucket $A$: P2 chooses $-50 \implies \text{utility} = -50$.
    \item Bucket $B$: P2 chooses $1 \implies \text{utility} = 1$.
    \item Bucket $C$: P2 chooses $-5 \implies \text{utility} = -5$.
\end{itemize}
Player 1 chooses bucket $B$ because its guaranteed minimum ($1$) is higher than $A$'s ($-50$) and $C$'s ($-5$). P2 subsequently selects $1$, resulting in final utility $1$.

\section{Tic-Tac-Toe}

Both players play autonomously using Minimax / Alpha-Beta search.

\subsection{Output}
\begin{lstlisting}[style=terminaloutput]
   |   |  
---+---+---
   |   |  
---+---+---
   |   |  

It is P1's turn to move
P1's action: (0, 0)

 x |   |  
---+---+---
   |   |  
---+---+---
   |   |  

It is P2's turn to move
P2's action: (1, 1)

 x |   |  
---+---+---
   | o |  
---+---+---
   |   |  

It is P1's turn to move
P1's action: (0, 1)

 x | x |  
---+---+---
   | o |  
---+---+---
   |   |  

It is P2's turn to move
P2's action: (0, 2)

 x | x | o
---+---+---
   | o |  
---+---+---
   |   |  

It is P1's turn to move
P1's action: (2, 0)

 x | x | o
---+---+---
   | o |  
---+---+---
 x |   |  

It is P2's turn to move
P2's action: (1, 0)

 x | x | o
---+---+---
 o | o |  
---+---+---
 x |   |  

It is P1's turn to move
P1's action: (1, 2)

 x | x | o
---+---+---
 o | o | x
---+---+---
 x |   |  

It is P2's turn to move
P2's action: (2, 1)

 x | x | o
---+---+---
 o | o | x
---+---+---
 x | o |  

It is P1's turn to move
P1's action: (2, 2)

 x | x | o
---+---+---
 o | o | x
---+---+---
 x | o | x

The game is a draw
\end{lstlisting}

\subsection{Discussion}
Under optimal play from both sides, neither player can force a win. The game correctly concludes in a draw after filling the board. Both Minimax and Alpha-Beta produce this identical sequence of moves.

\section{Non-Mandatory: Immediate Win in Tic-Tac-Toe}

Consider the board state:
\begin{center}
\begin{tabular}{c|c|c}
\texttt{x} & \texttt{o} & \texttt{o} \\ \hline
\texttt{x} &           &           \\ \hline
           &           &          
\end{tabular}
\end{center}

\subsection{Why Minimax Plays the Middle Square (1, 1)}
Standard utility assigns $+1$ to any terminal win, regardless of how many moves it takes. An immediate win at $(2, 0)$ gives utility $+1$, but playing the center $(1, 1)$ also guarantees a forced win later, which also evaluates to $+1$. Because both actions yield identical utility $+1$, Minimax treats them as tied and simply picks the first one encountered in row-major order: $(1, 1)$ before $(2, 0)$.

\subsection{Fixing the Game Class}
To prioritize faster wins, we reward wins that leave more empty squares (shorter depth):

\begin{lstlisting}[language=Python]
class FastWinGame(Game):
    def utility(self, state: State, player: int) -> float:
        assert self.is_terminal(state)
        empty = sum(row.count(None) for row in state[1])
        if self.is_winner(state, player):
            return 1 + empty
        if self.is_winner(state, (player + 1) % 2):
            return -(1 + empty)
        return 0
\end{lstlisting}

\subsection{Output Comparison}
With the standard \texttt{Game} class, P1 delays the win by playing $(1, 1)$ first:
\begin{lstlisting}[style=terminaloutput]
 x | o | o
---+---+---
 x |   |  
---+---+---
   |   |  

It is P1's turn to move
P1's action: (1, 1)

 x | o | o
---+---+---
 x | x |  
---+---+---
   |   |  

It is P2's turn to move
P2's action: (1, 2)

 x | o | o
---+---+---
 x | x | o
---+---+---
   |   |  

It is P1's turn to move
P1's action: (2, 0)

 x | o | o
---+---+---
 x | x | o
---+---+---
 x |   |  

P1 won
\end{lstlisting}

With \texttt{FastWinGame}, P1 plays $(2, 0)$ immediately to win on the first turn (utility $1 + 4 = 5$):
\begin{lstlisting}[style=terminaloutput]
 x | o | o
---+---+---
 x |   |  
---+---+---
   |   |  

It is P1's turn to move
P1's action: (2, 0)

 x | o | o
---+---+---
 x |   |  
---+---+---
 x |   |  

P1 won
\end{lstlisting}

\section{Runtime: Minimax vs. Alpha-Beta Pruning}

\subsection{First Move Runtimes}
Timing the search for the very first move on an empty $3 \times 3$ board:

\begin{table}[htbp]
\centering
\begin{tabular}{lccc}
\toprule
\textbf{Algorithm} & \textbf{First Move} & \textbf{Time (s)} & \textbf{Speedup} \\
\midrule
Minimax            & $(0, 0)$            & $15.34\text{ s}$  & $1.0\times$ \\
Alpha-Beta Pruning & $(0, 0)$            & $0.42\text{ s}$   & $\mathbf{36.1\times}$ \\
\bottomrule
\end{tabular}
\end{table}

\subsection{Why the Speedup Is So Drastic}
\begin{enumerate}
    \item \textbf{Subtree elimination:} Alpha-Beta does not just skip single moves; whenever $v \ge \beta$ or $v \le \alpha$, it skips entire recursive subtrees containing thousands of board states.
    \item \textbf{Branching factor reduction:} Alpha-Beta reduces the effective search complexity from $O(b^d)$ towards $O(b^{d/2})$ under good move ordering.
    \item \textbf{Frequent cutoffs from draws:} In Tic-tac-toe, many lines lead to a draw ($0$). Once $\alpha = 0$, any MIN branch finding a drawing move immediately triggers an $\alpha$-cutoff ($v \le \alpha$), skipping all other opponent replies.
    \item \textbf{Less state copying:} Skipping hundreds of thousands of node visits avoids calls to \texttt{deepcopy(board)}, significantly cutting Python runtime overhead.
\end{enumerate}

\end{document}
